\documentclass[12pt,a4paper,oneside]{report}

% FONT I JEZIK
\usepackage[utf8]{inputenc}
\usepackage[T1]{fontenc}
\usepackage{mathptmx}

\usepackage[croatian]{babel}
\usepackage{microtype}


% MARGINE  (izvorni dokument: 2,5 cm sa svih strana za glavni tekst;
%           naslovnice imaju 3 cm - to je postavljeno posebno kod naslovnica)
\usepackage{geometry}
\geometry{
  a4paper,
  top=2.5cm,
  bottom=2.5cm,
  left=2.5cm,
  right=2.5cm,
  headheight=14.5pt,
  headsep=1cm,
  footskip=1.1cm
}

% PRORED I ULOMCI (TEKST stil: prored 1,5, bez uvlake, mali razmak prije)
\usepackage{setspace}
\onehalfspacing
\usepackage[skip=3pt, indent=0pt]{parskip}
\emergencystretch=3em

% ZAGLAVLJE / PODNOŽJE  (Ime Prezime | Završni rad ... Fakultet ... | broj str.)
\usepackage{fancyhdr}
\pagestyle{fancy}
\fancyhf{}
\fancyhead[L]{\textit{Petra Alimović}}
\fancyhead[R]{\textit{Završni rad}}
\fancyfoot[L]{\textit{Fakultet strojarstva i brodogradnje}}
\fancyfoot[R]{\textit{\thepage}}
\renewcommand{\headrulewidth}{0.5pt}
\renewcommand{\footrulewidth}{0.5pt}


\fancypagestyle{plain}{%
  \fancyhf{}
  \fancyhead[L]{\textit{Petra Alimović}}
  \fancyhead[R]{\textit{Završni rad}}
  \fancyfoot[L]{\textit{Fakultet strojarstva i brodogradnje}}
  \fancyfoot[R]{\textit{\thepage}}
  \renewcommand{\headrulewidth}{0.5pt}
  \renewcommand{\footrulewidth}{0.5pt}
}

% NASLOVI POGLAVLJA / PODPOGLAVLJA
\usepackage{titlesec}

\titleformat{\chapter}[hang]
  {\bfseries\fontsize{14}{16}\selectfont}{\thechapter}{0.5em}{}
\titlespacing*{\chapter}{0pt}{30pt}{18pt}

\titleformat{\section}[hang]
  {\normalsize\bfseries}{\thesection}{0.5em}{}
\titlespacing*{\section}{0pt}{18pt}{12pt}

\titleformat{\subsection}[hang]
  {\normalsize\bfseries\itshape}{\thesubsection}{0.5em}{}
\titlespacing*{\subsection}{0pt}{12pt}{12pt}

\titleformat{\subsubsection}[hang]
  {\normalsize\itshape}{\thesubsubsection}{0.5em}{}
\titlespacing*{\subsubsection}{0pt}{12pt}{12pt}

\setcounter{secnumdepth}{3}
\setcounter{tocdepth}{3}

\renewcommand{\thechapter}{\arabic{chapter}.}
\renewcommand{\thesection}{\arabic{chapter}.\arabic{section}.}
\renewcommand{\thesubsection}{\arabic{chapter}.\arabic{section}.\arabic{subsection}.}
\renewcommand{\thesubsubsection}{\arabic{chapter}.\arabic{section}.\arabic{subsection}.\arabic{subsubsection}.}

% NEPREKIDNO NUMERIRANJE SLIKA / TABLICA / JEDNADŽBI
\usepackage{chngcntr}
\counterwithout{figure}{chapter}
\counterwithout{table}{chapter}
\counterwithout{equation}{chapter}

% NASLOVI SLIKA / TABLICA  ("Slika 1.  Kako opisati", podebljano, 11pt, centr.)
\usepackage{caption}
\captionsetup{
  labelsep=period,
  justification=centering,
  font={small,bf},
  labelfont={bf}
}
\captionsetup[figure]{name=Slika}
\captionsetup[table]{name=Tablica}

% OSTALI PAKETI
\usepackage{graphicx}
\usepackage{amsmath}
\usepackage{amssymb}
\usepackage{enumitem}
\usepackage{longtable}
\usepackage{array}
\usepackage{tocloft}
\usepackage[hidelinks]{hyperref}
\usepackage{bookmark}
\usepackage{minted}
\usepackage{xcolor}
\setminted{
    fontsize=\footnotesize,
    linenos=true,          % line numbers
    breaklines=true,       % wrap long lines
    frame=lines,           % horizontal rules above and below
    framesep=2mm,
    baselinestretch=1.2,
    style=friendly,              % VSCode-like light theme — see note below
}

\renewcommand{\cftchapdotsep}{\cftdotsep}
\renewcommand{\cftchapfont}{\normalfont}
\renewcommand{\cftchappagefont}{\normalfont}


% POMOĆNE NAREDBE
\newcommand{\naslovbeznumeracije}[1]{%
  \chapter*{#1}
  \addcontentsline{toc}{chapter}{#1}
}

\begin{document}


\renewcommand{\contentsname}{SADRŽAJ}
\renewcommand{\listfigurename}{POPIS SLIKA}
\renewcommand{\listtablename}{POPIS TABLICA}
\renewcommand{\bibname}{LITERATURA}

% 1) NASLOVNICA (naslovna stranica)  -- margine 3 cm, bez zaglavlja/podnožja
\newgeometry{top=3cm,bottom=3cm,left=3cm,right=3cm}
\pagestyle{empty}
\pagenumbering{gobble}

\begin{center}
  {\fontsize{16}{19}\selectfont SVEUČILIŠTE U ZAGREBU\\[2pt]
   FAKULTET STROJARSTVA I BRODOGRADNJE}
\end{center}

\vspace{2.4cm}

\begin{center}
  {\bfseries\fontsize{28}{34}\selectfont ZAVRŠNI RAD}
\end{center}

\vspace{2cm}

\begin{flushright}
  {\bfseries\fontsize{14}{17}\selectfont Petra Alimović}
\end{flushright}

\vfill

\begin{center}
  {\fontsize{14}{17}\selectfont Zagreb, godina 2026.}
\end{center}

\clearpage

% 2) NASLOVNICA (unutarnja naslovna stranica s mentorima)
\begin{center}
  {\fontsize{16}{19}\selectfont SVEUČILIŠTE U ZAGREBU\\[2pt]
   FAKULTET STROJARSTVA I BRODOGRADNJE}
\end{center}

\vspace{2.4cm}

\begin{center}
  {\bfseries\fontsize{28}{34}\selectfont ZAVRŠNI RAD}
\end{center}

\vspace{2cm}

\noindent
\begin{minipage}[t]{0.5\textwidth}
  {\fontsize{14}{17}\selectfont Mentori:}\\[6pt]
  {\fontsize{14}{17}\selectfont
   Prof.\ dr.\ sc.\ Marko Švaco, dipl.\ ing.}
\end{minipage}%
\begin{minipage}[t]{0.5\textwidth}\raggedleft
  {\fontsize{14}{17}\selectfont Student:}\\[6pt]
  {\fontsize{14}{17}\selectfont Petra Alimović}
\end{minipage}

\vfill

\begin{center}
  {\fontsize{14}{17}\selectfont Zagreb, godina 2026.}
\end{center}

\clearpage
\restoregeometry

% 3) IZJAVA
\pagestyle{empty}
\setlength{\parindent}{0.5cm}

\vspace*{2.2cm}

Izjavljujem da sam ovaj rad izradio samostalno koristeći znanja stečena
tijekom studija i navedenu literaturu.

\vspace{6pt}
Zahvaljujem se mentoru prof.\ dr.\ sc.\ Marku Švaci što mi je pružio
priliku za izradu ovog rada, te na njegovom vodstvu, uvidima i savjetima.

Zahvaljujem se laboratoriju CRTA i svim njegovim zaposlenicima
na pomoći i opremi potrebnoj za izradu ovog rada.

Zahvaljujem se svojoj obitelji na dugogodišnjoj podršci tijekom pisanja ovog rada i tijekom cijelog studija.

Naposljetku, posebno se zahvaljujem svim prijateljima i kolegama koji su mi pomogli svojim idejama, savjetima i podrškom.
Bez njih, ni početak niti završetak ovog rada ne bi bio moguć.

\vspace{1.2cm}
\begin{flushright}
Petra Alimović
\end{flushright}

\setlength{\parindent}{0pt}
\clearpage

% 4) ZAVRŠNI ZADATAK

\vspace*{\fill}
\begin{center}
\fbox{\begin{minipage}{0.85\textwidth}
\centering
\vspace{1cm}
\textit{[Ovdje umetnuti "Završni zadatak"]}
\vspace{1cm}
\end{minipage}}
\end{center}
\vspace*{\fill}
\clearpage

% 5) PRETHODNI SADRŽAJ (rimska paginacija: I, II, III ...)
\pagestyle{fancy}
\pagenumbering{Roman}
\setcounter{page}{1}

% --- SADRŽAJ ---------------------------------------------------------------
\phantomsection
\addcontentsline{toc}{chapter}{\contentsname}
\tableofcontents
\clearpage

% --- POPIS SLIKA -------------------------------------------------------------
\phantomsection
\addcontentsline{toc}{chapter}{\listfigurename}
\listoffigures
\clearpage

% --- POPIS TABLICA -----------------------------------------------------------
\phantomsection
\addcontentsline{toc}{chapter}{\listtablename}
\listoftables
\clearpage

% --- POPIS OZNAKA ------------------------------------------------------------
\naslovbeznumeracije{POPIS OZNAKA}
\begin{center}
\begin{tabular}{@{}p{3cm}p{3cm}l@{}}
\textbf{Oznaka} & \textbf{Jedinica} & \textbf{Opis} \\[4pt]
Oznaka 1 & Jedinica & Opis oznake \\
Oznaka 2 & Jedinica & Opis oznake \\
Oznaka 3 & Jedinica & Opis oznake \\
Oznaka 4 & Jedinica & Opis oznake \\
\end{tabular}
\end{center}
\clearpage

% --- SAŽETAK -----------------------------------------------------------------
\naslovbeznumeracije{SAŽETAK}
Tekst sažetka.

\vspace{6pt}
\noindent\textbf{Ključne riječi:}
\clearpage

% --- SUMMARY -----------------------------------------------------------------
\naslovbeznumeracije{SUMMARY}
Abstract.

\vspace{6pt}
\noindent\textbf{Key words:}
\clearpage

% 6) GLAVNI TEKST
\pagenumbering{arabic}
\setcounter{page}{1}

\chapter{UVOD}
Mobilni roboti su automatizirani, programabilni,
multifunkcionalni mehatroni\-čki sustavi sposobni kretati se
autonomno, to jest bez pomoći ljudskih operatera, kroz 
prostor kako bi izvršavali zadane zadatke.
[citati iz literature] 
To znači da robot sam ima sposobnost određivanja 
postupaka potrebnih za izvršavanje zadatka, 
koristeći senzore i algoritme za percepciju okoline. [citati iz literature]

Kada se ne kreću u potpunosti autonomno, njima je potrebno upravljati 
koristeći upravljačke palice, tipkovnicu, miša ili sličnu opremu. 
U novije vrijeme, mobilni roboti su sve pristupačniji te je njihova 
primjena sve šira i uključuje područja poput industrije, 
medicine, vojske, poljoprivrede, ali i u kućanstvima te za osobnu uporabu. 
To stvara potrebu za razvojem novih metoda upravljanja koja 
će omogućiti jednostavnije i intuitivnije upravljanje mobilnim robotima 
uz korištenje što manje opreme. Iz tog razloga, razvijeni su načini 
upravljanja koji se oslanjaju na glasovne naredbe, pokrete glave i očiju te 
geste ruku i tijela. 
Upravljanje mobilnim robotima pomoću gesti ruku posebno je zanimljivo 
jer omogućuje korisniku da upravlja robotom bez potrebe za 
složenom opremom, a pritom zadržava visoku razinu preciznosti. 


Problem kod upravljanja ASTRO robotom je što ne postoji ugrađeni sustav za zaštitu od sudara 
pri nailasku na prepreku. Ukoliko korisnik, ili sam robot, 
napravi krivi pokret, krivo procijeni udaljenost ili brzinu kretanja, 
nema dovoljno dobru povratnu informaciju o okolini i slično, 
može doći do zabijanja u prepreke, oštećenja robota, ozljeda korisnika ili druge štete. 
Kako bi se to izbjeglo, potrebno je razviti sustav za zaobilaženje prepreka. 
Tehnike za izbjegavanje prepreka baziraju se na detekciji prepreka u prostoru 
i planiranju putanje kretanja koristeći podatke dobivene sa senzora 
robota kako bi se izbjegao sudar. Implementacijom ovog sustava, 
mobilni robot dobiva veći stupanj autonomije.
Ovakav sustav primjenjiv je na sve robote, pa tako i mobilni robot ASTRO, koji imaju 
implementiran LiDAR senzor, neovisno o njihovom stupnju autonomije. 
Sustav djeluje kao zaštitni sloj (engl. safety layer) visoke razine prioriteta koji 
prema potrebi preuzima kontrolu nad kretanjem robota, zanemarujući ručne naredbe operatera 
dok se ne otkloni opasnost od sudara.

Cilj ovog završnog rada jest razviti i implementirati sustave za upravljanje mobilnim robotom 
u stvarnom vremenu koristeći upravljačke naredbe dobivene gestama ruke te za 
detekciju, automatsko zaobilaženje prepreka i izbjegavanje sudara.

\pagebreak

\chapter{MOBILNI ROBOTI}
Mobilni roboti se od ostalih razlikuju svojom sposobnošću kretanja po prostoru 
koja im je omogućena integracijom pogonskih, senzorskih i upravljačkih sustava. 
Svaki od tih sustava međusobnom integracijom omogućuje pravilan rad robota, a nalazi 
se u konstrukciji robota koja ga štiti od vanjskih utjecaja.

Pogonski sustav odgovoran je za svako gibanje robota. Aktuatori koji se najčesće koriste 
su elektromotori istosmjerne sutruje, sa ili bez četkica. Na njih su većinom 
ugrađeni enkoderi kojima se mjeri kut zakreta motora, a time i brzina, akceleracija te procjena pozicije.

Glavna funkcija senzorskog sustava je prikupljanje informacija iz okoline koje 
obrađuje upravljački sustav. Senzori su od velike važnosti jer omogućuju robotu da 
bude svjestan svoje okoline i promjena u njoj, što mu je potrebno za autonomno donošenje 
odluka. Najčešće korišteni senzori su LiDAR, kamere, ultrazvučni senzori, infracrveni senzori,
IMU (engl. Inertial Measurement Unit) i GPS (engl. Global Positioning System).

Upravljački sustav povezuje dva prethodno navedena sustava. Njegova osnova je mikrokontroler 
koji obrađuje podatke dobivene sa senzora i na temelju njih donosi odluke te upravlja aktuatorima 
kako bi ih proveo. Ovisno o složenosti robota, upravljački sustav može biti implementiran na mikrokontroleru,
računalu ili na više njih. U slučaju kada se koristi više mikrokontrolera ili računala, sustav dijelimo na 
upravljanje visoke razine (engl. high-level control) i upravljanje niske razine (engl. low-level control). 
Najčešće korišteni prograski jezici za programiranje upravljačkih sustava su C++ i Python.

Glavna podjela mobilnih robota vrši se na temelju medija u kojem se koriste, tlo, zrak ili voda. 
Podvodni i zračni roboti kreću se u tri dimenzije, dok su kopneni ograničeni na planarno kretanje. 
Zbog široke primjenjivosti, kopneni mobilni roboti su najrašireniji i najčešće korišteni, pa 
tako i u ovom radu.


\begin{figure}[htbp]
  \centering
  \fbox{\parbox[c][4cm]{6cm}{\centering Ovdje umetnuti sliku\\
  (\texttt{\textbackslash includegraphics})}}
  \caption{Slika}
  \label{fig:slika1}
\end{figure}

\section{ASTRO robot}
ASTRO robot (Autonomus System for Teaching RObotics) je mobilni robot razvijen u regionalnom 
centru izvrsnosti za robotske tehnologije CRTA za potrebe edukacije i istraživačkog razvoja. 
Za svrhu ovog rada, ASTRO nudi robusnu platformu jednostavnu za korištenje na koju su 
implementirani sustav za upravljanje gestama ruke i sustav za izbjegavanje prepreka.

\section{Arhitektura ASTRO robota}
Nadređeno upravljačko računalo na robotu je jednopločno računalo Nvidia Jetson Nano s operacijskim 
sustavom Ubuntu Linux i ROS 2 (Robot Operating System 2) Humble, sposobno za cijelokupnu obradu podataka.
Jetson Nano je malo računalo napravljeno za učenje, razvoj alata i proizvoda te drugih primjena u granama 
umjetne inteligencije i robotike. 

Upravljanje niske razine u stvarnom vremenu obavlja se putem mikrokontrolera Teensy 4.0, izgrađenog oko procesora 
ARM Cortex-M7 i spojenog na Jetson Nano putem sučelja UART (Universal Asynchronous Reciever Transmitter). 
Teensy pretvara zadanje linearne i kutne brzine lijevog i desnog kotača te zatvara PiD pretlju regulacije 
brzine na temelju povratne informacije s enkodera.

Cijela arhitektura ASTRO robota prikazana je na Slika [broj].

Za pokretanje ROS 2 čvorova koji omogućuju serijsku komunikaciju između 
komponenti koristi se Docker. Logika nadređenog sustava implementirana je u ROS 2 prostoru, 
a pomoću pokrenutih čvorova primaju se informacije sa senzora, upravljačkih palica ili informacije 
koje šalju vanjski spojeni uređaji. 

Prikaz cijele ROS 2 strukture vidi se na Slika [broj].

\chapter{SENZORSKI SUSTAV}
ASTRO koristi sljedeće senzore:
\begin{enumerate}
  \item enkoderi integrirani u pogonske motore
  \item Intel RealSense D435 Depth kamera
  \item Slamtec RPLIDAR A1 LiDAR.
\end{enumerate}
Svaki od ovih senzora je korišten tijekom izrade ovog rada. 

Osim senzora integriranih na ASTRO-u, korištena je još jedna Intel RealSense D435 Depth kamera postavljena 
izvan robota i spojena na vanjsko računalo koje je povezano s robotom i upravlja njime.

\section{Enkoderi}
Enkoderi mjere rotaciju svakog pogonskog kotača, iz čega se izračunava i akumulira prijeđeni put radi 
određivanja odometrije robota. Svaki motor opremljen je magnetskim enkoderom koji daje 16 impulsa po okretaju 
uz reduktor s prijenosnim omjerom 43,8:1, a prijeđeni putevi pojedinih kotača komniniraju se pomoću modela 
diferencijalnog pogona kako ni se izračunala pozicija robora. Enkoderi pružaju visoku brzinu osvježavanja te su 
jednostavni za integraciju. U ovom radu, odometrija s enkodera koristila se pri simulaciji kretanja robota putem 
razvijenog sustava za upravljanjem pomoću gesti ruke.

\section{Intel RealSense D435 Depth kamera}
Intel RealSense D435 širokokutna je stereo kamera za percepciju dubine, namijenjena primjenama u robotici 
i računalnom vidu. Kombinira par infracrvenih stereo kamera s aktivnim infracrvenim projektorom i RGB kamerom 
u boji, što joj omogućuje generiranje gustih, metričkih mapa dubine pri brzini do 90 slika u sekundi i 
dosegu od otprilike 0,2 do 10 metara. Dubina se izračunava na samom uređaju pomoću namjenskog čipa za 
sterokorekciju, koji na scenu projicira strukturirani infracrveni uzorak kako bi poboljšao teksturu i 
preciznost procjene dubine u područjima bez izražene teksture. Dobiveni trodimenzionalni oblak točaka, 
zajedno s poravnatom slikom u boji, dostupan je spojenom računalu putem Intel RealSense SDK 2.0 i 
njegovih Python biblioteka (pyrealsense2), koje su u ovom radu korištene za prikupljanje sinkroniziranih 
okvira boje i dubine. Kamera se na računalo spaja putem USB 3.0 sučelja i ne zahtijeva vanjsko 
napajanje, što je čini praktičnim i prenosivim senzorom za mobilne robotske platforme. U ovom radu korištene su 
dvije ovakve kamere. Prva kamera nalazila se na ASTRO robotu te služila za povratnu informaciju o poziciji robota 
u prostoru koja je dostupna operateru unutar korisničkog sučelja. Druga kamera bila je postavljena izvan 
robota i korištena za snimanje operaterove ruke u stvarnom vremenu, pri čemu je okvir dubine pružao 
metričku trodimenzionalnu poziciju ruke u odnosu na kameru, koja je potom korištena za izračun upravljačkih 
naredbi za robota.

\section{LiDAR senzor}
Slamtec RPLIDAR A1 je dvodimenzionalni laserski daljinomjer koji je sposoban izvesti potpuni sken okoline u 
rasponu od 360 stupnjeva. Radi na principu optičke triangulacije, tj. emitira infracrvenu lasersku zraku i 
mjeri kut pod kojim se reflektirana svjetlost vraća na pozicijski osjetljivi detektor, iz čega se izračunava 
udaljenost do reflektirajuće površine. Senzor se neprekidno rotira brzinom od otprilike 5,5 okretaja u 
sekundi, prikupljajući do 8000 mjerenja udaljenosti u sekundi pri kutnoj rezoluciji od približno jednog 
stupnja, s radnim rasponom od 0,15 do 12 metara. Sirovi podaci skena prenose se putem UART sučelja i 
dostupni su u ROS 2 okruženju kao standardna poruka tipa 	sensor\_msgs/LaserScan posredstvom službenog 
upravljačkog paketa 	sllidar\_ros2 koji pruža Slamtec, generirajući temu koju je izravno koristio čvor za 
izbjegavanje prepreka razvijen u ovom radu. Dobiveni sken pruža planarni presjek okoline na visini montaže 
senzora, što je dovoljno za detekciju prepreka i mjerenje slobodnih udaljenosti u horizontalnoj ravnini. 
U ovom radu RPLIDAR A1 bio je montiran na gornju stranu robota ASTRO i korišten kao primarni izvor za 
detekciju prepreka, omogućujući algoritmu za izbjegavanje praćenje okoline robota u svim smjerovima i 
određivanje sigurnih putanja kretanja.

\chapter{ROS2 SUSTAV}
ROS 2 (Robot Operating System 2) skup je alata i biblioteka otvorenog koda namijenjen izgradnji i 
integraciji aplikacija za robote. Unatoč svom nazivu, ROS 2 nije operativni sustav u tradicionalnom smislu, 
već skup alata, biblioteka i komunikacijskog međusloja (engl. middleware) koji neovisnim softverskim modulima, zvanim čvorovi 
(engl. nodes), omogućuje razmjenu podataka na strukturiran i standardiziran način. Svaki čvor je samostalan 
izvršni proces odgovoran za jedan, jasno definirani zadatak. Primjerice, u ovom su radu to su namjenski čvorovi 
zaduženi za snimanje gesta ruke, generiranje upravljačkih naredbi brzine, izbjegavanje prepreka te 
multipleksiranje izvora brzine. 

Primarni mehanizam komunikacije u ROS 2 su teme (engl. topics), koje 
implementiraju asinkroni model objavi-pretplati (engl. publish-subscribe). Čvor koji proizvodi podatke 
objavljuje poruke na imenovanom kanalu teme, dok se bilo koji broj drugih čvorova može pretplatiti na taj 
isti kanal i primati te poruke bez ikakve izravne sprege između objavitelja i pretplatnika. Ovakva odvojena 
arhitektura sustav čini inherentno modularnim: pojedini čvorovi mogu se razvijati, testirati, mijenjati ili 
iznova pokretati neovisno, bez narušavanja ostatka sustava. 

Svi podaci koji se razmjenjuju putem tema 
prenose se u porukama (engl. messages). One su strogo tipizirane podatkovne strukture definirane jezikom za 
definiranje sučelja ROS-a (IDL, engl. Interface Definition Language). Vrsta poruke definira polja i njihove 
tipove podataka, osiguravajući da svi čvorovi koji dijele neku temu budu suglasni oko formata razmijenjenih 
podataka. Primjeri poruka korištenih u ovom radu uključuju geometry\_msgs/Twist za upravljačke naredbe brzine, 
sensor\_msgs/LaserScan za podatke o dometima LiDAR senzora te prilagođenu vrstu poruke HandPose koja nosi 
trodimenzionalnu poziciju operaterove ruke zajedno s prepoznatom gestom i ocjenom pouzdanosti. 

Uz teme, ROS 2 pruža i servise (engl. services), koji implementiraju sinhroni komunikacijski obrazac 
zahtjev-odgovor: jedan čvor djeluje kao poslužitelj i izlaže imenovanu servisnu krajnju točku, dok 
klijentski čvor šalje tipizirani zahtjev i čeka tipizirani odgovor prije nastavka izvođenja. Servisi su 
prikladni za rijetke, jednokratne operacije poput upita o konfiguraciji ili promjene načina rada, za razliku 
od kontinuiranih tokova podataka za koje su teme namijenjene. 

U ovom radu upravljanje naredbama gibanja robota preuzima paket twist\_mux, koji se pretplaćuje na nekoliko 
konkurentnih tema tipa geometry\_msgs/Twist koje generiraju različiti čvorovi: čvor za upravljanje joystickom, 
čvor za upravljanje gestama ruke i čvor za izbjegavanje prepreka. Paket prosljeđuje samo poruku izvora s 
najvišim konfiguriranim prioritetom upravljaču robota. Svaki izvor ima pridruženo konfigurirano vremensko 
ograničenje. Ako izvor prestane objavljivati unutar tog vremenskog okvira, automatski gubi prednost i 
upravljanje prelazi na sljedeći dostupni izvor, čime se osigurava sigurno ponašanje u slučaju kvara čvora 
ili prekida komunikacije. 

Verzija ROS 2 koja se koristi u ovom radu je Humble Hawksbill, objavljena u svibnju 2022. godine s 
dugotrajnom podrškom zajamčenom do svibnja 2027. godine.

\chapter{MEDIAPIPE BIBLIOTEKA}
MediaPipe je međuplatformski okvir otvorenog koda koji je razvio Google za izgradnju cjevovoda percepcije 
(engl. perception pipeline) u stvarnom vremenu temeljenih na strojnom učenju. Umjesto jedne namjenske 
aplikacije, nudi biblioteku unaprijed istreniranih modela i modularnih komponenti za obradu zasnovanih na 
grafovima, koje se mogu međusobno povezivati radi formiranja potpunih rješenja za zadatke poput detekcije 
lica, procjene položaja tijela, detekcije objekata i praćenja ruke. Ključni cilj dizajna MediaPipea je 
učinkovitost. Njegovi su modeli optimizirani za rad u stvarnom vremenu na standardnom procesorskom 
sklopovlju bez potrebe za grafičkim procesorom, što ga čini praktičnim za primjenu na ugradbenim i 
mobilnim platformama kakve su uobičajene u robotici.

Komponenta korištena u ovom radu je MediaPipe Hands, rješenje za praćenje ruke koje radi kroz dvostupanjski 
cjevovod strojnog učenja. U prvom stupnju, lagani model za detekciju dlana pronalazi ruku na ulaznoj slici i 
vraća područje omeđivanja. U drugom stupnju, zasebni model za detekciju ključnih točaka ruke uzima izrezanu 
i normaliziranu sliku detektiranog područja te regresijom određuje položaje 21 trodimenzionalne ključne 
točke koje opisuju skeletnu strukturu ruke. Te su ključne točke postavljene na anatomski definirane položaje: 
zapešće, metakarpofalangealni (MCP) zglobovi kostura dlana, proksimalni interfalangealni (PIP) zglobovi, 
distalni interfalangealni (DIP) zglobovi i vrhovi prstiju za svaki od četiri prsta, zajedno s četiri 
odgovarajuće točke duž palca. Koordinate x i y svake ključne točke izražene su u normaliziranom prostoru 
slike u rasponu od nule do jedan u odnosu na dimenzije slike, dok koordinata z kodira relativnu dubinu unutar 
ruke, pri čemu zapešće služi kao referentno ishodište.

Nakon što je ruka jednom detektirana, cjevovod prelazi iz načina detekcije u način praćenja za sljedeće 
sličice, koristeći prethodne položaje ključnih točaka za vođenje učinkovitije operacije izrezivanja i 
poravnavanja, umjesto ponovnog pokretanja punog detektora dlana u svakom kadru. Detekcija se nastavlja tek 
kada ruka napusti kadar ili kada pouzdanost modela padne ispod konfigurabilnog minimalnog praga. Klasifikator 
lateralnosti dodatno dodjeljuje oznaku lijeve ili desne ruke detektiranoj ruci zajedno s ocjenom pouzdanosti. 
U ovom su radu oba praga pouzdanosti, za detekciju i za praćenje, postavljena na 0,6, a sličice ispod te 
vrijednosti odbacivane su i tretirane kao kadrovi bez detektirane ruke, kako bi se spriječilo prenošenje 
pogrešnih klasifikacija gesta na robota.

Svih 21 ključnih točaka korišteno je u ovom radu u dvije svrhe: klasifikaciju gesta i procjenu položaja dlana. 
Za klasifikaciju gesta, četiri MCP ključne točke kostura dlana na indeksima 5, 9, 13 i 17 usrednjene su kako bi 
se definiralo središte dlana. Za svaki od četiri prsta, euklidska udaljenost od ključne točke vrha prsta do 
tog središta dlana uspoređivana je s udaljenošću od odgovarajućeg PIP zgloba do istog središta. Ako je vrh 
prsta bio bliže dlanu nego PIP zglob za konfigurabilni faktor skale, prst je smatran savijenim. Ruka u kojoj 
su tri ili više prstiju istovremeno bili savijeni klasificirana je kao šaka, koja je služila kao ulaz 
sigurnosne sklopke mrtvog čovjeka, upravljačke naredbe brzine slane su robotu isključivo dok je operater 
održavao zatvorenu šaku.

Za procjenu položaja dlana, budući da MediaPipe Hands pruža samo normalizirane dvodimenzionalne koordinate 
slike i relativnu procjenu unutarnje dubine, a ne metričke trodimenzionalne položaje, za dobivanje stvarnih 
udaljenosti korišten je poravnati okvir dubine Intel RealSense D435 kamere. Koordinate piksela svake ključne 
točke korištene su za pronalazak odgovarajuće metričke vrijednosti dubine u okviru dubine, koja je zatim 
transformirana u trodimenzionalni prostor pomoću unutarnje matrice parametara kamere putem funkcije 
rs2\_deproject\_pixel\_to\_point iz RealSense SDK-a. Dobiveni trodimenzionalni položaj središta dlana, 
izražen u metrima u odnosu na koordinatni sustav kamere, objavljivao se kao tema u ROS 2 okruženju i 
koristio ga je čvor za daljinsko upravljanje za izračun proporcionalnih upravljačkih naredbi brzine za robota.

\chapter{OPEN CV BIBLIOTEKA}
OpenCV (Open Source Computer Vision Library) je sveobuhvatna biblioteka otvorenog koda namijenjena obradi 
slika i računalnom vidu u stvarnom vremenu, koja pruža nekoliko tisuća optimiziranih algoritama za zadatke 
poput filtriranja slike, detekcije rubova, kalibracije kamere, optičkog toka i prepoznavanja objekata. 
Biblioteka je izvorno napisana u programskom jeziku C++, no dostupna je i putem Python sučelja (cv2) koje 
je korišteno u ovom radu. Temeljni podatkovni model OpenCV-a oslanja se na NumPy polja: svaki okvir slike 
predstavljen je kao trodimenzionalno polje oblika (visina, širina, broj kanala), pri čemu su vrijednosti 
piksela pohranjene kao 8-bitni cijeli brojevi bez predznaka. Važna konvencija biblioteke je da boje interno 
pohranjuje redoslijedom plavog, zelenog i crvenog kanala (BGR), za razliku od uobičajenog RGB redoslijeda 
koji koriste druge biblioteke poput MediaPipea; zbog toga je u ovom radu svaki okvir boje prije 
prosljeđivanja modelu za detekciju ključnih točaka ruke pretvoren iz BGR u RGB prostor boja pozivom 
funkcije cv2.cvtColor s parametrom cv2.COLOR\_BGR2RGB.

Okviri dubine koje pruža Intel RealSense D435 pohranjena su kao 16-bitna polja s vrijednostima izraženima u 
milimetrima, što ih čini neprikladnima za izravni vizualni prikaz. U ovom je radu okvir dubine za potrebe 
grafičkog sučelja obrađen u dva koraka: funkcijom cv2.convertScaleAbs vrijednosti su linealno skalirane i 
svedene na 8-bitni raspon (0-255), nakon čega je funkcijom cv2.applyColorMap primijenjena pseudokolorna 
mapa\ cv2.COLORMAP\_JET, koja bliže udaljenosti prikazuje hladnim (plavim) bojama, a dalje udaljenosti 
toplim (crvenim) bojama. Na taj način vizualizirani okvir dubine prikazuje se uz okvir boje, dajući 
operateru intuitivan pregled prostorne geometrije scene u stvarnom vremenu.

Grafičko sučelje u ovom radu izgrađeno je isključivo korištenjem primitiva za crtanje OpenCV-a, bez 
upotrebe posebnog okvira za grafičko korisničko sučelje. Funkcija cv2.rectangle korištena je za crtanje 
obojenih pravokutnih traka kao pozadine za tekstualne informacije, čime su postignuti vizualno odvojeni 
informacijski paneli. Funkcija cv2.putText s fontom cv2.FONT\_HERSHEY\_SIMPLEX i zastavicom cv2.LINE\_AA za 
glađenje rubova korištena je za prikazivanje podataka o prepoznatoj gesti, pouzdanosti detekcije, statusu 
sigurnosne sklopke, odmaku ruke od neutralnog položaja i statusu kalibracije. Isti tekst iscrtavan je dva 
puta — jednom crnom bojom s većom debljinom linije kako bi se dobila sjena, a zatim u željenom obojenju s 
manjom debljinom — čime je postignut efekt obruba koji tekst čini čitkim neovisno o podlozi. Funkija 
cv2.circle korištena je za preslikavanje ključnih točaka ruke iz okvira boje na odgovarajuće položaje u 
koloriranom okviru dubine, čime se operateru vizualno potvrđuje da su detektirane točke prostorno 
konzistentne između oba prikaza. Konačno, dva vizualna panela — okvir boje s ucrtanim skeletom ruke i 
prikazanim informacijama te kolorirani okvir dubine — horizontalno su spojena pomoću NumPy funkcije 
np.hstack, dok je ispod njih vertikalnom kompozicijom (np.vstack) dodan statusni panel koji prikazuje 
kalibracijski napredak i odmake od neutralnog položaja u oba smjera. Rezultirajući kompozitni okvir 
prikazuje se pozivom cv2.imshow unutar prozora kreiranog s cv2.namedWindow, a tipkovnički unosi operatera 
obrađuju se funkcijom cv2.waitKey, koja je u ovom radu konfigurirana s vremenskim ograničenjem od jedne 
milisekunde kako ne bi blokirala petlju kamere — pritisak tipke C pokreće postupak kalibracije neutralnog 
položaja, dok tipka Q ili tipka Escape pokreće uredno gašenje čvora.

\chapter{SUSTAV ZA UPRAVLJANJE POMOĆU GESTI RUKE}
Problem upravljanja mobilnim robotom pomoću gesti ruke rješava se razvojem dva ROS 2 čvora koji koriste 
MediaPipe Hands za detekciju ruke i klasifikaciju gesta, te Intel RealSense D435 Depth kameru za procjenu 
trodimenzionalnog položaja dlana. 

Prije razvijanja i pisanja koda za čvorove, osmišljena je logika sustava, definirane su geste koje će se 
koristiti za upravljanje robotom, teme i poruke koje će se koristiti za komunikaciju između čvorova te 
napisan pseudokod koji opisuje način rada svakog čvora. Također, instalirani su svi potrebni ROS 2 paketi 
i biblioteke, uključujući MediaPipe i Intel RealSense SDK 2.0, te su testirani njihovi osnovni primjeri
kako bi se osiguralo da su ispravno instalirani i funkcionalni. 

Prvi razvijeni čvor, hand\_gesture\_recognizer, odgovoran je za prepoznavanje gesta ruke i objavljivanje 
rezultata u ROS 2 okruženju. Drugi čvor, hand\_gesture\_controller, pretvara prepoznate geste i 
procijenjeni položaj dlana u upravljčke naredbe brzine za robota, koje se objavljuju na standardnoj 
ROS 2 temi geometry\_msgs/Twist.

Prije implementacije sustava na fizičkog mobilnog robota ASTRO, razvijeni su i testirani u simuliranom
okruženju pomoću rviz-a, gdje su se provjeravale ispravnost prepoznavanja gesta, točnost procjene 
položaja dlana i odgovarajuće generiranje upravljačkih naredbi brzine. Nakon što je sustav uspješno testiran 
u simulaciji, implementiran je na fizički robot ASTRO, gdje je dodatno testiran u stvarnom vremenu s operaterom.

Zbog jednostavnosti, čvorovi će u ovom radu biti opisani pojedinačno i preko pseudokoda. Poveznica za cijeli 
paket i kodove svih čvorova nalazi se u Prilogu I.

\section{Prepoznavanje geste ruke}
Logika osmišljena za prepoznavanje gesta ruke temelji se na klasifikaciji ruke u dva stanja: 
otvorena i zatvorena šaka. Kao i većina robotskih sustava, i ovaj sustav koristi sigurnosnu sklopku 
mrtvog čovjeka (engl. dead man's switch), što znači da se poruke o položaju dlana, a time i brzine robota, 
šalju isključivo dok operater drži šaku zatvorenom. Ukoliko se šaka otvori, robot se odmah zaustavlja i ne 
prima nove upravljačke naredbe dok se šaka ponovno ne zatvori. 

Osim geste šake, sustav prati i udaljenost ruke od kalibriranog početnog položaja u horizontalnoj ravnini. 
Iz te udaljenosti se kasnije izračunava proporcionalna brzina kretanja robota, što je detaljnije opisano u 
poglavlju 6.2.

U ovom čvoru je korišten i OpenCV za obradu slike i vizualni prikaz povratnih informacija operateru preko 
grafičkog korisničkog sučelja. OpenCV je korišten za prikaz slike iz kamere, označavanje detektirane ruke i
prikaz prepoznate geste. Također, korišten je i za prikaz poravnate slike dubine te video prijenos s kamere 
pozicionirane s prednje strane robota kako bi operater imao bolju povratnu informaciju o poziciji robota u 
prostoru.

Uvezeni moduli u ovom čvoru su:
\begin{itemize}
  \item rclpy
  \item rclpy.node
  \item geometry\_msgs.msg.Point
  \item std\_msgs.msg.String
  \item std\_msgs.msg.Header
  \item sensor\_msgs.msg.CompressedImage
  \item pyrealsense2
  \item mediapipe
  \item cv2
  \item numpy
  \item time
\end{itemize}

Nakon inicijalizacije čvora, definirane su sljedeće varijable i parametri:

\begin{minted}{python}
self.PUBLISH_RATE           = 20
self.MIN_CONFIDENCE         = 0.6
self.DEBOUNCE_FRAMES        = 5
self.FIST_CURL_THRESHOLD    = 0.8
self.CALIBRATION_DURATION   = 3.0

self.confirmed_gesture  = 'OPEN'
self.pending_gesture    = 'OPEN'
self.pending_count      = 0
self.latest_position    = None
self.latest_confidence  = 0.0
self.last_known_position = Point()

self.null_position          = None
self.is_calibrating         = False
self.calibration_start_time = 0.0
\end{minted}

Jedni od važnijih parametara koji nisu striktno za kameru 
su MIN\_CONFIDENCE, DEBOUNCE\_FRAMES i FIST\_CURL\_THRESHOLD. MIN\_CONFIDENCE definira 
minimalnu pouzdanost klasifikacije geste ruke kako bi se smatrala valjanom. To smanjuje broj pogrešnih 
klasifikacija koje se mogu pojaviti zbog lošeg osvjetljenja, preklapanja ruke s drugim objektima ili 
nepreciznosti u detekciji. Parametar je određen eksperimentalno, ukoliko je pouzdanost pre visoka, sustav 
nikada neće detektirati zatvorenu šaku. Isto tako, ukoliko je pouzdanost pre malena, sustav će lako krivo 
detektirati gestu koju operater ne pokazuje, što može dovesti do nesigurnog ponašanja robota i trzaja. 
DEBOUNCE\_FRAMES definira broj uzastopnih 
sličica koje moraju biti klasificirane kao šaka ili otvorena ruka prije nego li se promijeni stanje potvrđene 
geste. Taj parametar je važan jer MediaPipe Hands može ponekad pogrešno klasificirati ruku, a DEBOUNCE\_FRAMES 
omogućuje da se pogreške filtriraju, što sprečava trzaje robota pri kretanju. FIST\_CURL\_THRESHOLD definira 
prag za klasifikaciju geste šake na temelju omjera udaljenosti vrha prsta i PIP zgloba od središta dlana. 

Ostali parametri definirani u ovom čvoru su za kalibraciju početnog položaja ruke, praćenje trenutne pozicije 
i pohranu posljednje poznate pozicije ruke.

Nakon definiranja početnih parametara, kreirani su ROS 2 pretplatnici i objavljivači.

\begin{minted}{python}
Create Publisher: '/hand_pose' (HandPose message)
Create Publisher: '/null_position' (Point message)
Create Subscriber: '/camera/color/image_raw/compressed' -> robot_camera_callback()
Create Subscriber: '/obstacle_status'                  -> obstacle_status_callback()
\end{minted}

HandPose poruka se objavljuje na temu /hand\_pose i sadrži trodimenzionalnu poziciju ruke, prepoznatu gestu 
i ocjenu pouzdanosti. Kasnije se koristi za izračun upravljačkih naredbi brzine robota. Null\_position poruka 
se objavljuje na temu /null\_position i sadrži trodimenzionalnu poziciju ruke u trenutku kada je operater 
započeo kalibraciju.
Pretplatnici se kreiraju na temu /camera/color/image\_raw/compressed, koja sadrži komprimirane slike iz
Intel RealSense D435 kamere, te na temu /obstacle\_status, koja sadrži status prepreke i koristi se za 
zaustavljanje robota u slučaju da sustav za izbjegavanje prepreka detektira prepreku u blizini robota. 

Dalje u kodu slijedi inicijalizacija MediaPipe Hands modela, intrinzika i poravnanja kamere te kreiranje tajmera.
\begin{minted}{python}
Initialize RealSense pipeline (Color & Depth @ 640x480, 30fps)
Initialize MediaPipe Hands model (min_detection_confidence, max_num_hands=1)
        
Create Timer(1/30s) -> camera_loop()
Create Timer(1/20s) -> publish_latest_state()
\end{minted}

Pri inicijalizacije MediaPipe Hands modela, postavljen je maksimalan broj ruku koji može biti detektiran u 
jednom trenutku. MediaPipe Hands model radi na principu da ukoliko detektira ruku, ona se prati dok ne 
napusti kadar ili dok pouzdanost klasifikacije ne padne ispod praga. Parametar je postavljen na 1 jer se u 
ovom radu koristi samo jedna ruka operatera te se na ovaj način
sprječava pogrešno detektiranje druge ruke operatera ili ruke druge osobe u kadru.

Glavna funkcija programa, camera\_loop(), poziva se svakih 1/30 sekundi i obrađuje sliku iz kamere.

\begin{minted}{python}
function camera_loop():
  color_frame, depth_frame = RealSense.get_aligned_frames()
  IF frames are invalid:
      Debounce_Gesture('OPEN')
      RETURN
\end{minted}

Prvi dio funkcije je akvizicija i poravnanje kadra. Funkcija RealSense.get\_aligned\_frames() dohvaća
okvire boje i dubine iz Intel RealSense D435 kamere i poravnava ih tako da svaki piksel u okviru boje 
odgovara istom pikselu u okviru dubine. Ukoliko su okviri neispravni, funkcija Debounce\_Gesture poziva 
se s argumentom 'OPEN' kako bi se osiguralo da robot prestane primati upravljačke naredbe i zatim se 
funkcija prekida i vraća. Također, potrebno je pretvoriti kadrove u numpy polja, napraviti konverziju 
iz BGR u RGB prostor boja i napraviti vizualizaciju okvira dubine za prikaz u grafičkom korisničkom sučelju.

Nakon toga, MediaPipe Hands model obrađuje okvir boje i vraća rezultate detekcije ruke i ključnih točaka. 

\begin{minted}{python}
  results = MediaPipe.process(color_frame)
  IF hand NOT detected OR confidence < MIN_CONFIDENCE:
      Debounce_Gesture('OPEN')
      Render_Dashboard_GUI()
      RETURN
\end{minted}

Ukoliko je pouzdanost detekcije ruke manja od praga MIN\_CONFIDENCE, funkcija Debounce\_Gesture poziva se s
argumentom 'OPEN' kako bi se osiguralo da robot prestane primati upravljačke naredbe. Također, potrebno je 
pozvati funkciju Render\_Dashboard\_GUI() kako bi se prikazali vizualni podaci u grafičkom korisničkom sučelju.

\begin{minted}{python}
  landmarks_2d = results.multi_hand_landmarks[0]
  raw_gesture  = Classify_Gesture(landmarks_2d)
  Debounce_Gesture(raw_gesture)
\end{minted}

Sljedeći korak je dohvat ključnih točaka ruke i klasifikacija geste. Funkcija Classify\_Gesture() koristi
omjer udaljenosti vrha prsta i PIP zgloba od središta dlana kako bi odredila je li ruka otvorena ili 
zatvorena šaka. 

\begin{minted}{python}
  function Classify_Gesture(landmarks_2d):
        palm_center = Calculate_Mean(landmarks_2d[5, 9, 13, 17])
        curled_fingers_count = 0

        FOR EACH finger IN [index, middle, ring, pinky]:
            tip_dist = Distance(finger.tip, palm_center)
            pip_dist = Distance(finger.pip, palm_center)

            IF tip_dist < pip_dist * FIST_CURL_THRESHOLD:
                curled_fingers_count += 1

        IF curled_fingers_count >= 3:
            RETURN 'FIST'
        ELSE:
            RETURN 'OPEN'
\end{minted}

Šaka je definirana kao zatvorena ukoliko su istovremeno savijena tri ili više prstiju. 
Taj broj je određen eksperimentalno, kada je prag bio manji, sustav je pogrešno klasificirao ruku kao šaku 
te kada je prag bio veći, sustav je pogrešno klasificirao ruku kao otvorenu. Stroži prag bi bio bolji za 
sigurnost, no za to su potrebni bolji uvjeti okoline za kameru ili bolji senzori.
Nadalje, ako je jedan ili manje prstiju savijeno, ruka se klasificira kao otvorena, a sve ostalo se tretira kao neodređeno. 
Neodređeno stanje na upravljanje utjeće isto kao i otvorena ruka, tj. robot prestaje primati upravljačke naredbe. 
Nadalje, ponovno se poziva funkcija Debounce\_Gesture() kako bi se osiguralo da se pogrešne 
klasifikacije filtriraju i da se robot ne trza pri kretanju. 

\begin{minted}{python}
  landmarks_3d = Deproject_2D_to_3D(landmarks_2d, depth_frame, camera_intrinsics)
  palm_center  = Compute_Mean_3D_Position(landmarks_3d[knuckles])
  state.position_3d = palm_center
\end{minted}

Dalje, budući da MediaPipe Hands model pruža samo normalizirane dvodimenzionalne koordinate slike i 
relativnu procjenu unutarnje dubine, a ne metričke trodimenzionalne položaje, za dobivanje stvarnih 
udaljenosti koristi se poravnati okvir dubine Intel RealSense D435 kamere. Funkcija Deproject\_2D\_to\_3D() 
koristi koordinate piksela svake ključne točke kako bi pronašla odgovarajuću metričku vrijednost dubine u 
okviru dubine, koja se zatim transformira u trodimenzionalni prostor pomoću unutarnje matrice parametara 
kamere putem funkcije rs2\_deproject\_pixel\_to\_point iz RealSense SDK-a.
Nakon toga, izračunava se središte dlana kao srednja vrijednost 3D pozicija MCP zglobova 
kostura dlana, koje se nalaze na indeksima 5, 9, 13 i 17. Dobivena pozicija središta dlana pohranjuje se u 
varijablu state.position\_3d. 

\begin{minted}{python}
  IF calibration_is_active:
      IF elapsed_time >= CALIBRATION_DURATION:
          state.null_position = state.position_3d
          Publish('/null_position', state.null_position)
          Stop_Calibration()
\end{minted}

Kako bi se iz pozicije ruke izračunale upravljačke naredbe brzine, potrebno je odrediti referentnu poziciju 
ruke, tj. neutralni položaj, od kojeg se računa udaljenost ruke u horizontalnoj ravnini. 
Funkcija za kalibraciju aktivira se pritiskom tipke C na tipkovnici, nakon čega se mjeri vrijeme određeno 
parametrom CALIBRATION\_DURATION. Nakon što je proteklo dovoljno vremena, trenutna pozicija ruke se pohranjuje 
u varijablu state.null\_position i objavljuje na temu /null\_position. Nakon toga, kalibracija se zaustavlja 
i sustav je spreman za objavljivanje informacija o poziciji i gesti ruke. 

\begin{minted}{python}
  Render_Dashboard_Grid_2x2(color_frame, depth_colormap, robot_image, obstacle_status)
      
  key = Read_Keyboard_Input()
  IF key == 'C': Start_Calibration_Timer()
  IF key == 'Q': Shutdown_Node()
\end{minted}

Kako bi sustav bio što intuitivniji i jednostavniji za korištenje, razvijeno je grafičko korisničko sučelje 
koje prikazuje četiri panela. 
Gornji lijevi panel prikazuje video prijenos s vanjske kamere s ucrtanim skeletom ruke i ključnim točkama te 
tekstualne informacije o prepoznatoj gesti, pouzdanosti detekcije i upute za operatera poput unosa s tipkovnice 
za kalibraciju i gašenje sustava. 
Gornji desni panel prikazuje kolorirani okvir dubine s ucrtanim ključnim točkama ruke, što operateru daje 
informaciju o prostornoj poziciji ruke u odnosu na okolinu. Uz kolorirani okvir dubine nalazi se i legenda 
koja prikazuje povezanost boja i udaljenosti. 
Donji lijevi panel prikazuje video prijenos s kamere pozicionirane s prednje strane robota, što operateru 
daje bolju povratnu informaciju o poziciji robota u prostoru. Također, služi operateru za nadođenje u situacijama 
kada se robot ne nalazi u operaterovom vidnom polju. U tom slučaju, taj video prijenos je jedini način da 
operater orijentira robota u točnom smjeru. 
Donji desni panel jest statusni panel koji prikazuje razne informacije o stanju sustava poput statusa kalibracije i 
statusa sustava za detekciju prepreka. U ovom panelu operater dobiva informaciju kada mu se oduzima mogućnost upravljanja 
robotom u slučaju nailaska na prepreku, u kojem je stadiju algoritam za zaobilaženje prepreke te, naposlijetku, 
kada mu je vraćena kontrola nad robotom. 

Naposlijetku, definirana je funkcija koja objavljuje najnovije stanje sustava.

\begin{minted}{python}
  function publish_latest_state():
        msg = HandPose()
        msg.timestamp = Current_Time()
        msg.gesture   = state.confirmed_gesture
        msg.confidence = state.latest_confidence
        msg.position  = state.position_3d IF valid ELSE state.last_known_position

        Publish('/hand_pose', msg)
\end{minted}

Ovaj čvor prikuplja, obrađuje i objavljuje sve potrebne informacije za računanje brzina robota i upravljanje 
mobilnim robotom ASTRO pomoću pozicije i geste ruke. 

\section{Izračun brzina robota}
Jednom kada je određena udaljenost šake od referentnog položaja u planarnoj ravnini i je li šaka otvorena 
ili zatvorena, iz tih podataka može se izračunati brzina robota, linearna i kutna, proporcionalna određenoj 
udaljenosti. Sve koordinate pozicije šake koje se koriste u ovom čvoru izražene su u koordinatnom sustavu 
optičkog okvira kamere (engl. \textit{camera optical frame}), koji je standardna konvencija za RGB-D kamere. 
U tom koordinatnom sustavu os x pokazuje udesno, os y prema dolje, a os z u dubinu scene, tj. od kamere 
prema promatranoj sceni. Iz toga slijedi da pomak šake u smjeru osi x odgovara lateralnom pomaku (lijevo-desno), 
koji se preslikava na kutnu brzinu robota, dok pomak u smjeru osi z odgovara promjeni udaljenosti šake od 
kamere, koji se preslikava na linearnu brzinu robota.

Uvezeni moduli u ovom čvoru su:
\begin{itemize}
  \item rclpy
  \item rclpy.node
  \item geometry\_msgs.msg.Twist
  \item geometry\_msgs.msg.Point
  \item hand\_teleop\_msgs.msg.HandPose
  \item math
\end{itemize}

Modul math koristi se za matematičke funkcije korištene kod računanja brzine.

Nakon inicijalizacije čvora, potrebno je odrediti početne parametre. Svi parametri su određeni eksperimentalno. 

\begin{minted}{python}
  class HandTeleopNode(ROS2_Node):

    function __init__():
        Initialize ROS 2 node "hand_teleop_node"

        Declare and Load Parameters:
            - deadzone_radius, max_linear_speed, max_angular_speed
            - map_range_linear, map_range_angular
            - smoothing_alpha (EMA filter coefficient)
            - tracking_timeout, publish_rate (20 Hz)

        state = {
            'neutral_position': NULL,
            'latest_hand_pos': NULL,
            'latest_gesture': "OPEN",
            'last_msg_time': NULL,
            'filtered_linear_x': 0.0,
            'filtered_angular_z': 0.0
        }
\end{minted}

Parametar deadzone\_radius potreban je kako bi se mali pomaci šake oko kalibrirane točke ignorirali. To je 
važno jer se ovaj način upravljanja bazira na gruboj motorici operatera te moramo filtrirati i ignorirati 
male drhtaje i pokrete ruku i tijela kako se robot ne bi trzao. 
Parametri \texttt{map\_range\_linear} i \texttt{map\_range\_angular} definiraju fizikalni raspon kretanja 
šake koji odgovara punoj brzini robota. Konkretno, \texttt{map\_range\_linear} predstavlja udaljenost u 
metrima kojom operater mora pomaknuti šaku prema ili od kamere u odnosu na referentni položaj kako bi robot 
dostigao svoju maksimalnu linearnu brzinu. Primjerice, postavljanjem vrijednosti na 0,20 m postiže se da 
pomak šake od 20 cm od referentne točke prema kameri rezultira maksimalnom brzinom kretanja robota naprijed. 
Isti princip vrijedi za \texttt{map\_range\_angular} u lateralnom smjeru. Manji raspon daje osjetljivije, 
brže upravljanje, dok veći raspon daje precizniju, finiju kontrolu pri manjim pomacima.
Smanjenju trzaja robota pridonosi i smoothing\_alpha parametar koji se koristi u smoothing filteru. 
Parametri za brzinu koriste se za skaliranje brzine u odnosu na udaljenost. Publish\_rate odgovara onome 
definiranom u prethodnom čvoru. 

Također su definirana i početna stanja tako da je pri pokretanju sustava sve u mirovanju radi sigurnosnih razloga. 

\begin{minted}{python}
        Create Subscriber: '/hand_pose'     -> on_hand_pose_received()
        Create Subscriber: '/null_position' -> on_null_position()
        Create Publisher:  '/cmd_vel_hand' (Twist message)

        Create Timer(1 / publish_rate) -> publish_velocity()
\end{minted}

Potom je potrebno stvoriti komunikaciju s ostalim čvorovima kako bi se dobile potrebne informacije. 
Čvor se pretplaćuje na /hand\_pose i /null\_position za informacije o gesti ruke i udaljenosti od referentnog 
položaja, isto kao i informaciju o referentnom položaju. 
Ovaj čvor objavljuje poruku oblika Twist na temu /cmd\_vel\_hand, na koju se pretplaćuje paket twist\_mux. 
Taj paket u svakom trenutku prati više konkurentnih izvora naredbi brzine poput upravljanje upravljačkom palicom, 
upravljanje gestama ruke i izbjegavanje prepreka. Paket prosljeđuje naredbu izvora s najvišim 
konfiguriranim prioritetom na temu /cmd\_vel, koju čita upravljač robota ASTRO. Na taj način, ukoliko 
operater istovremeno koristi upravljačku palicu i upravljanje gestama, uprednost ima način upravljanja s većim 
prioritetom. 
Uz pretplatnike i objavljivače, stvoren je i tajmer kako bi se čvor mogao sinkronizirati s hand\_capture čvorom. 

\begin{minted}{python}
  function on_null_position(msg):
        state.neutral_position = msg.position


    function on_hand_pose_received(msg):
        state.last_msg_time   = Current_Time()
        state.latest_hand_pos = msg.position
        state.latest_gesture  = msg.gesture
\end{minted}

Funkcije on\_null\_position i on\_hand\_pose\_recieved služe da dohvaćanje informacija dobivenih od tema na koje 
je čvor pretplaćen kako bi se mogle koristiti u daljnjem kodu. 

Potom se pokreće glavna funkcija programa, publish\_velocity, koja služi za objavu brzina motora robota. 
Funkcija se sastoji od sedam dijelova. 

\begin{minted}{python}
  function publish_velocity():
        IF state.neutral_position is NULL OR state.last_msg_time is NULL:
            Publish_Zero_Twist()
            RETURN

        IF (Current_Time() - state.last_msg_time) > tracking_timeout:
            Publish_Zero_Twist()
            RETURN

        IF state.latest_gesture != "FIST":
            Publish_Zero_Twist()
            RETURN
\end{minted}

Prvi dio funkcije ima sigurnosnu ulogu. Razmatraju se tri uvjeta pri kojima se umjesto stvarne brzine 
objavljuje nulta brzina. Prvi uvjet provjerava je li kalibracija referentnog položaja uopće provedena te 
je li pristigla barem jedna poruka s pozicijom šake. Ukoliko nije, robot miruje. Drugi uvjet jest vremensko 
ograničenje praćenja (\texttt{tracking\_timeout}), koje se aktivira kada od posljednje primljene poruke s 
pozicijom šake protekne više vremena nego što je dozvoljeno. Do toga može doći ako šaka napusti vidno polje 
kamere, ako MediaPipe izgubi praćenje ruke uslijed lošeg osvjetljenja ili brzog pokreta, ili ako dođe do 
prekida prijenosa podataka između čvorova. Vrijednost timeouta namjerno je postavljena malo ispod timeouta 
koji koristi \texttt{twist\_mux} (0,5 s) kako bi se osiguralo da ovaj čvor sam zaustavi robota i objavi 
nultu brzinu prije nego što \texttt{twist\_mux} samostalno reagira na nedostatak poruka. Treći uvjet jest 
sigurnosna sklopka mrtvog čovjeka. Robot prima naredbe za kretanje isključivo dok operater drži zatvorenu 
šaku. Otvaranje šake, prijelazno stanje geste ili gubitak detekcije rezultiraju trenutnom objavom nulte 
brzine.

\begin{minted}{python}
        dx = state.latest_hand_pos.x - state.neutral_position.x
        dz = state.neutral_position.z - state.latest_hand_pos.z
\end{minted}

Sljedeći dio funkcije služi za računanje pomaka, udaljenosti, od kalibrirane neutralne pozicije. Važno je napomenuti 
zašto se pomak u dubinu dz računa kao razlika u obrnutom redoslijedu u odnosu na dx. Kamera Intel RealSense 
D435 mjeri dubinu kao udaljenost od objektiva kamere, što znači da se vrijednost koordinate z povećava što 
je ruka dalje od kamere. Kako bi se postiglo intuitivno upravljanje, tj. pomicanje šake prema kameri treba 
rezultirati kretanjem robota naprijed, potrebno je preokrenuti predznak. Oduzimanjem trenutne pozicije šake 
od referentne vrijednosti (\texttt{neutral\_position.z - latest\_hand\_pos.z}) dobiva se pozitivan \texttt{dz} 
kada se šaka pomiče prema kameri, što se zatim preslikava na pozitivnu linearnu brzinu, odnosno kretanje 
robota prema naprijed. U suprotnom, pomicanje šake od kamere daje negativan \texttt{dz} i robot se kreće 
unatrag.

\begin{minted}{python}
        IF Euclidean_Distance(dx, dz) < deadzone_radius:
            dx = 0.0
            dz = 0.0
\end{minted}

Prije računanja brzine iz pomaka definiranog u prijašnjem koraku, potrebno je filtrirati pomake manje od 
radijusa mrtve zone koja je definirana u početnim parametrima. Ukoliko je pomak manji od zadanog radijusa, 
pomak će se smatrati da je jednak nuli.

\begin{minted}{python}
        norm_forward = Clamp(dz / map_range_linear, -1.0, 1.0)
        norm_turn    = Clamp(dx / map_range_angular, -1.0, 1.0)
\end{minted}

Potom je potrebno provesti normalizaciju i ograničavanje kojima se fizikalni pomak ruke u metrima pretvara 
u siguarn upravljački signal od -1.0 do 1.0. Normalizacija pretvara izmjerenu fizikalnu udaljenost ruke od nulte 
točke u relativni postotak u odnosu na zadani maksimalni raspon pokreta definiran početnim parametrima. 
Ograničavanje osigurava da dobivena normalizirana vrijednost strogom matematičkom granicom ne može izaći izvan 
raspona [-1.0, 1.0], bez obzira koliko daleko operater pomakne ruku izvan predviđenog radnog prostora. 
Svrha ove operacije je sigurnost i stabilnost robota, linearizacije pruža operateru proporcionalan 
osjećaj vožnjea sigurnosni strop sprječava da robot zbog prevelikog pokreta ruke dobije naredbu za brzinu 
koja premašuje maksimalne fizičke granice robota.

\begin{minted}{python}
        target_linear_x  = norm_forward * max_linear_speed
        target_angular_z = norm_turn * max_angular_speed
\end{minted}

Peti korak jest skaliranje normaliziranih vrijednosti na stvarne fizikalne brzine množenjem s parametrima 
\texttt{max\_linear\_speed} i \texttt{max\_angular\_speed}. Ove granične vrijednosti odgovaraju hardverskim 
ograničenjima robota ASTRO definiranim u upravljaču \texttt{astro\_hardware\_interface}. Maksimalna linearna 
brzina iznosi ±1,0 m/s, a maksimalna kutna brzina ±1,0 rad/s. Budući da je ulazna normalizirana vrijednost 
ograničena na raspon [-1.0, 1.0] u prethodnom koraku, umnožak s ovim graničnim vrijednostima jamči da 
izlazna naredba brzine nikada neće premašiti fizikalne granice robota, neovisno o položaju ruke operatera.

\begin{minted}{python}
        state.filtered_linear_x  = (smoothing_alpha * target_linear_x) + ((1 - smoothing_alpha) * state.filtered_linear_x)
        state.filtered_angular_z = (smoothing_alpha * target_angular_z) + ((1 - smoothing_alpha) * state.filtered_angular_z)
\end{minted}

Nakon skaliranja primjenjuje se eksponencijalno klizno usrednjavanje 
(EMA, engl. \textit{Exponential Moving Average}), opisano formulom:

\begin{equation}
y_k = \alpha \cdot x_k + (1 - \alpha) \cdot y_{k-1}
\end{equation}

gdje je xk trenutno izračunata ciljna brzina, yk-1 prethodno filtrirana brzina, a a koeficijent 
izglađivanja u rasponu (0,1]. Vrijednost a bliža nuli daje sporiji, glađi odziv filtera jer se više uzima 
u obzir prethodna vrijednost, dok vrijednost bliža jedan daje brži odziv, ali manje izglađen signal. Svrha 
ovog filtera jest smanjenje naglih skokova brzine koji bi mogli nastati zbog šuma u mjerenjima kamere ili 
brzih nevoljnih pokreta ruke operatera, čime se postiže glađe i sigurnije kretanje robota.

\begin{minted}{python}
        msg = Twist()
        msg.linear.x  = state.filtered_linear_x
        msg.angular.z = state.filtered_angular_z
        Publish('/cmd_vel_hand', msg)
\end{minted}

Završni korak funkcije jest objavljivanje brzina izračunatih ovom funkcijom kako bi ju robot mogao iskoristiti.

\begin{minted}{python}
  function Publish_Zero_Twist():
        state.filtered_linear_x  = 0.0
        state.filtered_angular_z = 0.0

        msg = Twist()
        msg.linear.x  = 0.0
        msg.angular.z = 0.0
        Publish('/cmd_vel_hand', msg)
\end{minted}

Naposlijetku, potrebno je vratiti memoriju filtera u početno stanje kako bi se prevenirao iznenadni skok 
akceleracije pri ponovnom pokretanju. 

Opisanim čvorom zaokružena je izrada cjelokupnog lanca obrade podataka koji omogućuje upravljanje 
robotom ASTRO pokretima ruke, od snimanja geste kamerom, preko klasifikacije i kalibracije, do izračuna i 
objave proporcionalnih naredbi brzine. Svi opisani čvorovi zajedno tvore funkcionalan sustav koji je 
spreman za integraciju u postojeću arhitekturu robota putem paketa \texttt{twist\_mux}. 

Radi lakšeg pokretanja sustava pomoću samo jedne naredbe, stvoren je paket koji sadrži dva čvora opisana u 
ovom poglavlju te datoteka za pokretanje.

\begin{minted}{python}
  def generate_launch_description():

    twist_mux_config = os.path.join(
        get_package_share_directory('astro'),
        'config',
        'twist_mux.yaml'
    )

    return LaunchDescription([

        Node(
            package='hand_teleop',
            executable='hand_capture_node',
            name='hand_capture_node',
            output='screen',
        ),

        Node(
            package='hand_teleop',
            executable='hand_teleop_node',
            name='hand_teleop_node',
            output='screen',
            parameters=[{
                'deadzone_radius'   : 0.03,
                'max_linear_speed'  : 1.0,
                'max_angular_speed' : 1.0,
                'map_range_linear'  : 0.20,
                'map_range_angular' : 0.20,
                'smoothing_alpha'   : 0.3,
                'tracking_timeout'  : 0.4,
                'publish_rate'      : 20.0,
            }]
        ),

        Node(
            package='twist_mux',
            executable='twist_mux',
            name='twist_mux',
            output='screen',
            parameters=[twist_mux_config],
            remappings=[('/cmd_vel_out', '/cmd_vel')],
        ),
    ])
\end{minted}

Pokretanjem ove datoteke pokreću se sljedeći čvorovi:
\begin{itemize}
  \item hand\_capture\_node
  \item hand\_teleop\_node
  \item twist\_mux.
\end{itemize}


\section{Simulacija}
Međutim, prije testiranja na stvarnom robotu, potrebno je provjeriti ispravnost implementiranog sustava u 
kontroliranim uvjetima kako bi se otklonili potencijalni problemi bez rizika od oštećenja opreme ili 
ugrožavanja sigurnosti. U tu svrhu provedena je provjera funkcionalnosti sustava u simulacijskom okruženju 
RViz, alatom za vizualizaciju tema i stanja robota unutar ROS 2 ekosustava. Budući da RViz za pravilno 
funkcioniranje zahtijeva prisutnost odometrijskih podataka koje stvarni robot generira putem enkodera, a 
koji nisu dostupni u simulacijskom okruženju bez pokretanja cjelokupnog hardverskog sloja, bilo je potrebno 
implementirati dodatni čvor koji generira sintetičke odometrijske podatke. 

\begin{minted}{python}
  class DummyOdomNode(ROS2_Node):

    function __init__():
        Initialize ROS 2 node "dummy_odom_node"

        state = {
            'x': 0.0, 'y': 0.0, 'th': 0.0,
            'vx': 0.0, 'wz': 0.0
        }

        last_time = Current_Time()

        Create Subscriber: '/cmd_vel_hand' (Twist message) -> cmd_vel_callback()
        Initialize TransformBroadcaster (tf2_ros)

        Create Timer(0.05s) -> update_and_publish()
\end{minted}

Početak programa inicijalizira sam čvor, određuje početne parametre, postavlja pretplatnika za dobivanje 
informacije o brzini, odašiljač transformacija i tajmer. Pozicija je zapisana u varijablama x i y, dok je 
orijentacija pohranjena u varijabli th. Linearna brzina je označena varijablom vx, a kutna varijablom wz.

\begin{minted}{python}
  function cmd_vel_callback(msg):
        state.vx = msg.linear.x
        state.wz = msg.angular.z
\end{minted}

Stvoren je pretplatnički povratni poziv koji se pokreče automatski svaki put kada čvor upravljanja pošalje 
novu brzinu na temu /cmd\_vel\_hand. Funkcija prihvaća poruku msg, izvuče brzine vx i wz te ih sprema u memoriju. 

\begin{minted}{python}
  function update_and_publish():
        current_time = Current_Time()
        dt = current_time - last_time
        last_time = current_time

        delta_x  = state.vx * Cosine(state.th) * dt
        delta_y  = state.vx * Sine(state.th) * dt
        delta_th = state.wz * dt

        state.x  += delta_x
        state.y  += delta_y
        state.th += delta_th
\end{minted}

U glavnoj funkciji, update\_and\_publish, prvo se računa proteklo vrijeme. Zatim se provodi kinematska 
integracija te se preko diferencijalnog pogonskog kinematičkog modela računaju brzine. 
Kinematska integracija koristi diferencijalnu jednadžbu gibanja diferencijalnog pogona. Na temelju trenutne 
orijentacije robota fi i izmjerene linearne brzine vx te kutne brzine wz, pomaci po osima izračunavaju se kao:

\begin{align}
\Delta x &= v_x \cdot \cos(\theta) \cdot \Delta t \
\Delta y &= v_x \cdot \sin(\theta) \cdot \Delta t \
\Delta \theta &= \omega_z \cdot \Delta t
\end{align}

gdje je dt proteklo vrijeme od prethodnog poziva funkcije. Integriranjem ovih diferencijalnih pomaka u 
svakom vremenskom koraku dobiva se procjena trenutne pozicije i orijentacije robota u ravnini, što je 
standardni postupak estimacije odometrije za mobilne robote s diferencijalnim pogonom.

\begin{minted}{python}
        transform = TransformStamped()
        transform.header.stamp    = current_time
        transform.header.frame_id = 'odom'
        transform.child_frame_id  = 'base_footprint'

        transform.translation = (state.x, state.y, 0.0)

        transform.rotation.x = 0.0
        transform.rotation.y = 0.0
        transform.rotation.z = Sine(state.th / 2.0)
        transform.rotation.w = Cosine(state.th / 2.0)
\end{minted}

Sljedeći korak jest konstrukcija TF okvira transformacije. 
Koordinatni okvir \texttt{odom} predstavlja fiksnu referentnu točku u prostoru iz koje robot počinje 
kretanje i u kojoj se nalazi ishodište odometrije. Okvir \texttt{base\_footprint} je koordinatni okvir 
pričvršćen za projekciju baze robota na tlo i pomiče se zajedno s robotom. Transformacija između ta dva 
okvira opisuje koliko se i u kojem smjeru robot pomaknuo od početne pozicije.
Potom slijedi pretvorba kuta zakretanja u kvaternion. Rotacija se u ROS 2 ne 
izražava Eulerovim kutovima već kvaternionima, četverokomponentnim matematičkim objektima koji jednoznačno 
opisuju orijentaciju u prostoru bez problema singularnosti karakterističnih za Eulerove kutove. Budući da 
se robot kreće samo u horizontalnoj ravnini, jedina relevantna rotacija je zakret oko osi z (kut fi), pa se 
kvaternion pojednostavljuje na oblik (0, 0, sin(fi/2), cos(fi/2)).

\begin{minted}{python}
        TF_Broadcaster.sendTransform(transform)
\end{minted}

Naposlijetku se odašilju koordinate transformacije ROS 2 TF stabla. ROS 2 TF stablo 
(engl. \textit{transform tree}) sustav je za praćenje prostornih odnosa između koordinatnih okvira robota 
u realnom vremenu. Svaki dio robota (baza, senzori, kotači) ima vlastiti koordinatni okvir, a TF stablo 
bilježi kako su ti okviri međusobno pozicionirani i orijentirani u svakom trenutku. RViz koristi TF stablo 
kako bi znao gdje se u prostoru nalaze pojedini dijelovi robota i kako ih vizualizirati. 
Bez odometrijske transformacije između okvira \texttt{odom} i \texttt{base\_footprint}, RViz ne može 
prikazati kretanje robota.


Na vanjsko računalo, na koje je bila spojena kamera Intel RealSense D435 i s kojeg se pokretao cijeli 
sustav, kloniran je repozitorij ASTRO s GitHub stranice CRTA laboratorija. Repozitorij sadrži sve potrebne 
datoteke za upravljanje robotom, uključujući i konfiguraciju za vizualizaciju u RViz okruženju s URDF 
modelom robota. Istovremeno su pokrenuti čvor \texttt{dummy\_odom}, datoteka \texttt{hand\_teleop.launch.py} 
koja pokreće čvorove za snimanje gesta i izračun brzina te \texttt{twist\_mux}, a u RViz okruženju učitan 
je model robota ASTRO. Ispravnost sustava potvrđena je vizualnim praćenjem modela robota u RViz prozoru. 
Model se kretao naprijed, nazad i zakretao u skladu s pokretima šake operatera, s vidljivom reakcijom na 
otvaranje šake (zaustavljanje) i zatvaranje šake (nastavak kretanja), čime je potvrđena funkcionalna 
ispravnost implementiranog lanca upravljanja prije testiranja na stvarnom robotu.

(Upute za pokretanje simulacije nalaze se na GitHub repozitoriju CRTA laboratorija. Uz RViz su pokrenute 
dummy\_odom.py i hand\_teleop.launch.py datoteke.) fusnota


\section{Integracija s mobilnim robotom ASTRO}
Budući da je simulacijom potvrđena ispravnost sustava, završni korak je bio integrirati sustav s fizičkim 
mobilnim robotom ASTRO. 
U ovom radu korišten je ASTRO 3 te će se u buduće koristiti taj identifikacijski (ID) broj, no ovaj sustav je 
primjenjiv na bilo koji mobilni robot ASTRO ukoliko se ID broj zamijeni valjanim.

U prethodnim poglavljima opisano je kako je kloniran ASTRO repozitorij te kako je izrađena datoteka za pokretanje. 
Stoga, jedini preostali korak jest pokrenuti ASTRO robota i bežično ga spojiti s vanjskim računalom. 

ASTRO se pokreće pritiskom tipki Power i JETSON. Tada je potrebno pričekati nekoliko 
minuta dok se ne pojavi Wi-Fi mreža pod nazivom Astro 3. Jednom kada je moguće, računalo se spaja na tu mrežu 
te je sklopovlje spremno za pokretanje sustava. 

Fizički robot ASTRO 3 ponašao se u skladu s očekivanjima potvrđenim simulacijom. Zatvaranjem dlana u šaku 
deaktivirala se sklopka mrtvog čovjeka te je robot počinjao primati naredbe za kretanje. Pomicanjem šake 
prema kameri robot se kretao naprijed, pomicanjem od kamere unatrag, a lateralnim pomicanjem šake ulijevo 
ili udesno robot se zakretao u odgovarajućem smjeru. Brzina kretanja bila je proporcionalna odmaku šake od 
referentnog položaja, a otvaranjem šake robot se trenutno zaustavljao, čime je potvrđena ispravna 
funkcionalnost sigurnosne sklopke mrtvog čovjeka u uvjetima stvarnog upravljanja. 

Optimalni parametri sustava određeni su eksperimentalno iterativnim postupkom testiranja na fizičkom robotu. 
Parametri koji su zahtijevali fino ugađanje uključuju \texttt{min\_confidence}, \texttt{debounce\_frames}, 
\texttt{fist\_curl\_treshold}, \texttt{deadzone\_radius}, \texttt{smoothing\_alpha} te \texttt{map\_range\_linear} i \texttt{map\_range\_angular}.
Konačne vrijednosti svih parametara rezultat su kompromisa između intuitivnosti upravljanja, 
stabilnosti kretanja i glatkoće vožnje robota.

Uspješnom integracijom i testiranjem na fizičkom robotu potvrđena je funkcionalna ispravnost cjelokupnog 
sustava upravljanja gestama ruke razvijenog u ovom radu. Sustav omogućuje intuitivno, bežično upravljanje 
mobilnim robotom ASTRO isključivo pokretima šake operatera, bez potrebe za fizičkim kontrolerom. 
Arhitektura temeljena na ROS 2 čvorovima i \texttt{twist\_mux} paketu osigurava modularnost i jednostavnu 
proširivost. Novi načini upravljanja ili novi senzori mogu se dodati bez izmjena postojećeg koda. 

zakljucak
(Razvijena rješenja, uključujući čvorove za snimanje i klasifikaciju gesta te izračun brzina i izbjegavanje prepreka, 
predstavljaju osnovu na kojoj se mogu graditi naprednije metode interakcije čovjeka i robota u budućim istraživanjima CRTA laboratorija.)
zakljucak

\chapter{SUSTAV ZA ZAOBILAŽENJE PREPREKA}
Objašnjen problem i rješenje.


\chapter{ZAKLJUČAK}
Tekst zaključka.

% ============================================================================
% 7) LITERATURA
% ============================================================================
\phantomsection
\addcontentsline{toc}{chapter}{\bibname}
\begin{thebibliography}{99}
\bibitem{kraut} Kraut, B.: Strojarski priručnik, Tehnička knjiga Zagreb, 1970.
\end{thebibliography}

% ============================================================================
% 8) PRILOZI
% ============================================================================
\naslovbeznumeracije{PRILOZI}
\begin{enumerate}[label=\Roman*.]
  \item link za git
\end{enumerate}

\end{document}
